\section{Full-Cost Floors and the Accountability Floor}\label{app:theory}

This section collects the results on the costs of running AI that Section~\ref{sec:horizons} of the main paper uses. Further supporting results and the notation are in Section~\ref{supp:theory}, and all proofs are in Section~\ref{app:proofs}.

\begin{lemma}[Components of the full-cost floors]\label{lem:fullcost}
All quantities are annual and per dollar of revenue.
\begin{enumerate}
\item[(a)] (Maintenance.) If automated builds cost $c_{\rm bld}$, need upkeep at the rate $\delta_M$ a year and a revalidation costing $\varrho_{\rm rv}c_{\rm bld}$ at each model upgrade, every $T_{\rm up}$ years, then $c_{\mathrm{mnt}}=(\delta_M+\varrho_{\rm rv}/T_{\rm up})\,c_{\rm bld}$; with discounting at $r$, $1/T_{\rm up}$ becomes $r/(e^{rT_{\rm up}}-1)$.
\item[(b)] (Security.) Let the expected incident loss without security spending be $c^{\rm inc}$, and let spending $c^{\rm sp}$ multiply every incident hazard by $\exp(-c^{\rm sp}/c_{\rm eff})$, with $c_{\rm eff}>0$. The risk cost $c^{\rm risk}:=\min_{c^{\rm sp}\ge0}\{c^{\rm sp}+e^{-c^{\rm sp}/c_{\rm eff}}c^{\rm inc}\}$ equals $c^{\rm inc}$ if $c^{\rm inc}\le c_{\rm eff}$ and $c_{\rm eff}[1+\ln(c^{\rm inc}/c_{\rm eff})]$ otherwise. Optimal spending is $c_{\rm eff}\ln(c^{\rm inc}/c_{\rm eff})^+\le c^{\rm inc}/e$, and when it is positive it leaves the residual loss $c_{\rm eff}$, whatever the exposure. If today's spending $c^{\rm sp}_0>0$ is optimal, so that $L_{\rm sec}:=\ln(c^{\rm inc}_0/c_{\rm eff})>0$, and full automation multiplies the expected loss by $\xi_{\rm exp}\ge1$, then $c^{\rm risk}=(c^{\rm sp}_0/L_{\rm sec})(1+L_{\rm sec}+\ln\xi_{\rm exp})$.
\item[(c)] (Errors and hallucination.) Let automated work whose in-house cost is $W_A$ carry, per unit, an expected $\varepsilon^{\rm o}$ ordinary errors and $\varepsilon^{\rm h}$ confident fabrications, and let each undetected error cost $L_{\rm err}$ times the unit's in-house cost. At full automation the review, if any, is by AI: $n_r\ge0$ review layers on one provider. On a share $\chi\in[0,1]$ of errors every review layer fails together with the reviewed model (it shares the model, its training data or its blind spot); on the other errors each layer catches an ordinary error with probability $d^{\rm o}$ and a fabrication with probability $d^{\rm h}$, independently. Whatever units are reviewed, the expected loss is at least $W_A\,\bar u\,L_{\rm err}$, where
\[
\bar u:=\varepsilon^{\rm o}\bigl[\chi+(1-\chi)(1-d^{\rm o})^{n_r}\bigr]+\varepsilon^{\rm h}\bigl[\chi+(1-\chi)(1-d^{\rm h})^{n_r}\bigr]
\]
is the review-proof error rate; $\bar u$ falls with depth to $\chi(\varepsilon^{\rm o}+\varepsilon^{\rm h})$, and with no review it is $\varepsilon^{\rm o}+\varepsilon^{\rm h}$. A requirement that undetected errors not exceed $\bar\varepsilon$ per unit can be met at depth $n_r$ if and only if $\bar u\le\bar\varepsilon$: some finite depth meets it if $\chi(\varepsilon^{\rm o}+\varepsilon^{\rm h})<\bar\varepsilon$, and none does if $\chi(\varepsilon^{\rm o}+\varepsilon^{\rm h})>\bar\varepsilon$.
\item[(d)] (Misbehavior.) Let each unit trigger a consequential misbehavior event with probability $\varepsilon^{\rm m}$, stopped before harm with probability $d^{\rm m}$, and otherwise costing $L_{\rm mis}Y$ times the unit's in-house cost, where $Y\ge1$ is Pareto with tail index $p_Y>0$, $P(Y>y)=y^{-p_Y}$, capped at $\bar Y>1$ by least privilege, independently. The expected loss is $W_A\,\varepsilon^{\rm m}(1-d^{\rm m})L_{\rm mis}\,\mathbb E\min\{Y,\bar Y\}$, with $\mathbb E\min\{Y,\bar Y\}=1+(1-\bar Y^{1-p_Y})/(p_Y-1)$ for $p_Y\ne1$ and $1+\ln\bar Y$ for $p_Y=1$. Without the cap the expectation is $p_Y/(p_Y-1)$ for $p_Y>1$ and infinite for $p_Y\le1$.
\item[(e)] (Common-mode failure.) Let the automated processes lose $W_{\rm cm}$ a year while all of them are disrupted, spread over model providers with shares $w_m$, and let $H_{\rm cm}:=\sum_mw_m^2$. Common-mode events arrive at each provider independently at the rate $\nu_{\rm cm}:=\nu_{\rm out}+(1-q_{\rm chg})\nu_{\rm reg}$: outages and shared vulnerabilities at the rate $\nu_{\rm out}$, and regressions after model updates at the rate $\nu_{\rm reg}$, of which change-management controls (pinned model versions, testing before release, rollback) stop the share $q_{\rm chg}\in[0,1]$ before they reach production. Each event disrupts every process on the provider for the containment time $T_{\rm con}$, and an event that disrupts the loss rate $x$ costs $T_{\rm con}[x+x^2/(2Y_{\rm fb})]$, where $Y_{\rm fb}$ measures fallback capacity. The expected annual cost is $\nu_{\rm cm}T_{\rm con}[W_{\rm cm}+H_{\rm cm}W_{\rm cm}^2/(2Y_{\rm fb})]$, which is largest with a single provider ($H_{\rm cm}=1$) and nonincreasing in $q_{\rm chg}$.
\end{enumerate}
\end{lemma}

\begin{remark}[The full-cost floors in numbers]\label{rem:fullcost-values}
Table~\ref{tab:fullcost} evaluates \autoref{lem:fullcost} for the cases of Table~\ref{tab:ceiling}. Values come from the parameter files of the replication package and from the assumptions register, whose grounding classes are shown: (a) sourced, (b) grounded by precedent, (c) calibrated from our own data, (d) a scenario value. Two rules set the three cases. Factors in classes (a) to (c) take the low end, the central value and the high end of their ranges. Scenario composites take a central value from the stated structure, halved and doubled. Fixed costs are evaluated at \$1 billion of revenue and fall as a share of revenue with scale. Two choices need comment. The conservative error rates are held at today's central values, because an activity at today's worst measured rates fails the test of \autoref{lem:fullcost}(c) and would not be automated. The risk cost uses \autoref{lem:fullcost}(b). In the conservative case, $L_{\rm sec}=1$, today's spending is at the bound of $1/e$ of the expected loss that \citet{gordon2002economics} derive for two classes of breach-probability functions. At full automation no human reviews the work, so the review-proof error rate of \autoref{lem:fullcost}(c) is that of AI review or of none. The catch rates that the literature measures, and that the replication package records, are those of human reviewers (spreadsheet inspection; crowdworkers and radiologists facing wrong AI suggestions), so they do not apply here. The ceiling therefore uses the rate with no review, $\bar u=\varepsilon^{\rm o}+\varepsilon^{\rm h}$, which needs no catch rate; AI review can lower it at most to $\chi(\varepsilon^{\rm o}+\varepsilon^{\rm h})$, and its compute belongs in $c_c$. For a review that flags the items on which two providers disagree, an error escapes with probability $P(\text{second model wrong}\mid\text{first wrong})\times P(\text{same wrong answer}\mid\text{both wrong})$, and only the second factor is measured, at about 0.60 on HELM \citep{kim2025correlated}. The provider-concentration term $H_{\rm cm}$ of \autoref{lem:fullcost}(e) is a lower bound on the effective concentration, because providers' errors are correlated and providers share infrastructure. The companion paper \citep{murthy2026transformation} sources each rate and derives the floors activity by activity.
\end{remark}

\begin{remark}[The accountability floor in numbers]\label{rem:amin-values}
The example follows the companion paper \citep{murthy2026transformation}; its values are illustrative, and it is not a control-design recommendation: admissibility is decided by management's assessment and the auditor's evidence. A firm with \$1 billion of revenue has one in-scope control objective with $N_o=24$ occasions a year on which a material misstatement arises (an illustrative count, for example two a month with the interim closes included), and $\pi$ is the probability that at least one of them passes every control layer in the year. Human control layers miss with probability 0.05 and cost \$500{,}000 a year each, with a common-mode probability of 0.002; AI layers on one provider miss with probability 0.02 and cost \$50{,}000, with a common-mode probability of 0.005 given their IT general controls, which cost \$204{,}672 a year, the line of one in-scope AI application priced at the end of this remark. There is no cross-type term, a human layer misses the items the AI layers missed at its own rate, and there is no separate cost of errors below materiality.

\emph{The threshold.} The standard sets no number for ``remote'', so $\bar\pi$ is a scenario value, and we report $\bar\pi=\AminEpsCen\%$ and $\bar\pi=\AminEpsSens\%$ a year side by side without choosing either as central. Two facts bear on the choice. First, adverse auditor attestations were 4.6\% of all SOX 404(b) attestations for fiscal 2020 \citep{auditanalytics2021sox404} and 5.8\% for fiscal 2021, against 4.8\% for 2019, in the latest edition we could read at the source \citep{auditanalytics2022sox404}; later editions exist, but we could not read their figures at the source and do not use them. These rates are an incidence of controls judged ineffective, not the probability that a misstatement goes undetected, so they bear on the plausibility of a threshold rather than measure it. In the example, every all-human design leaves at least $1-0.998^{24}\approx\AminHumLim\%$ a year, so none passes at \AminEpsSens\%, while three human layers pass at \AminEpsCen\% by a hair (\AminHumanProb\% a year, \$\AminHumanCost{} million). A threshold at which designs run by people fail would make ineffective control the norm among filers whose controls are mostly run by people, which the attestation rates contradict; at \AminEpsSens\% that is either a sign that the threshold is strict or that the human common-mode probability of 0.002 is high for an effective control environment. Second, the evidence is at the level of the firm and the constraint per objective: a firm-level \AminEpsCen\% spread evenly over \AminObjN{} objectives is \AminObjFive\% per objective (\AminObjOne\% at \AminEpsSens\%), at which the cheapest admissible design is \AminObjDesignFive{} (\$\AminObjCostFive{} million; \AminObjDesignOne{} and \$\AminObjCostOne{} million at \AminEpsSens\%).

\emph{Designs.} By \autoref{prop:amin}, an all-AI design on one provider is admissible only if $p^{\mathrm{cc}}+(1-p^{\mathrm{cc}})p^{\mathrm{cm}}<1-(1-\bar\pi)^{1/N_o}$, which is \AminPbarFive{} at \AminEpsCen\% and \AminPbarOne{} at \AminEpsSens\%. The example's common-mode probability of 0.005 exceeds both, so no all-AI design on the provider passes at either threshold, however many layers it uses: it lets at least one material misstatement through in a year with probability at least $1-0.995^{24}\approx\AminAiLim\%$. That conclusion is an assumption of the example, and the provider's common-mode probability under the firm's IT general controls is the number to estimate. At \AminEpsCen\% the cheapest admissible design is one human layer over one AI layer (\AminCenProb\% a year, \$\AminCenCost{} million, about \AminBest\% of revenue); at \AminEpsSens\% it is one human layer over two AI layers (\AminSensProb\% a year, \$\AminSensCost{} million, about \AminBestOne\%). With 48 occasions the design at \AminEpsCen\% needs a second AI layer (\$\AminAltNCost{} million); with an AI common-mode probability of 0.01 it still passes (\AminAltCProb\% a year). At \AminEpsSens\% either change requires two human layers over one AI layer (\$\AminAltOneCost{} million).

\emph{Break-evens.} Which design is cheapest is conditional, not a finding (Table~\ref{tab:amin-pcc}, whose $\bar\varepsilon$ is the threshold $\bar\pi$). At \AminEpsCen\%, one human over one AI layer stays admissible only while the cross-type common-mode probability is below about \AminPccOneOneFive{} per occasion and while the human layer's own miss rate on the items the AI layers missed, apart from its common mode of 0.002, is at most \AminPhmFive\% (\AminPhmTotFive\% including it). Automation bias, the tendency of a reviewer to accept wrong output from an automated aid, raises that miss rate. As the cross-type term rises, the cheapest design at \AminEpsCen\% becomes one human over two AI layers and then two humans over one AI layer, and once the cross-type probability reaches the threshold per occasion itself, \AminPbarFive{} at \AminEpsCen\% and \AminPbarOne{} at \AminEpsSens\%, no design is admissible. The cost of the cheapest design moves less than its identity.

\emph{Direct costs and the floor.} The expected-loss form of the floor in an earlier version of this paper charged \$10 million for each undetected misstatement and compared designs, two of which leave a reasonable possibility of a material misstatement; under the standard that loss does not enter the floor, and only the requirement does. Direct audit and Section 404 costs add about \AminDirect\% of revenue at this size: the average budget of the SOX programs of firms with less than \$5 billion of revenue in a survey of SOX programs, \$\KpmgSoxSmall{} million (values read from a text extraction of the archived report, whose chart layout and labels were verified; \citealp{kpmg2025sox}), plus the average audit fee of accelerated filers in fiscal 2022, \$\AuditFeeAccel{} million \citep{auditanalytics2023fees}. This fee is the low-side choice: a public software firm with \$1 billion of revenue is normally a large accelerated filer, whose average fee is the high case below. At the audit-fee ratio of the S\&P 500, \$343 per \$1 million of revenue in fiscal 2022 \citep{keohane2023sp500}, the direct cost is \AminDirectLo\%, and at the average fee of large accelerated filers, \$\AuditFeeLarge{} million, it is \AminDirectHi\%. The floor of one control objective is then about \AminFloor\% of revenue at \AminEpsCen\% and \AminFloorOne\% at \AminEpsSens\%, and \AminFloorLo{} to \AminFloorHi\% over that range of fees. The ceiling of Section~\ref{sec:horizons-ceiling} uses the band of $a_{\min}$, not this floor: set to the floor, the central threshold of Table~\ref{tab:ceiling} would be \GstarAminFloor\%.

An IT-general-controls line like that of an application performing invoice matching and exceptions costs about \ItgcLine\% of revenue a year. It covers access over the application's identities and tool permissions, change management, operations, four interfaces into in-scope systems, and a retest of its controls at each model upgrade, every 1.22 years. The retest costs \$\ItgcRetest: the 16 hours per control test that a survey of SOX programs reports on average \citep{kpmg2025sox}, applied to 10 controls at \$150 an hour; the number of controls and the hourly rate are our assumptions, and the other items of the line are scenario values. It covers control testing only, not the functional revalidation of the build counted in $c_{\mathrm{mnt}}$. A manual fallback for the close costs about \$161{,}000 more. The floor at \AminEpsCen\% reaches 0.5\% of revenue at about \ItgcApps{} in-scope AI applications on this line, \ItgcAppsFb{} with fallbacks, or \ItgcAppsRef{} on a line of \$80{,}000 a year (\ItgcLineRef\% of revenue), the survey-based reference of Table~\ref{tab:param-sources}; at the audit fee of a large accelerated filer it exceeds 0.5\% before any AI application. \texttt{code/verify\_theory.py} checks these values.
\end{remark}

\IfFileExists{sections/tables/tab_amin_pcc.tex}{\begin{table}[t]\centering\small
\caption{Cheapest admissible control design in the example of Remark~\ref{rem:amin-values}, by remote threshold, as the cross-type common mode $p^{\rm cc}$ and the human miss rate on AI-missed items vary}\label{tab:amin-pcc}
\begin{tabular}{@{}lll@{}}\toprule
 & $\bar\pi=1\%$ a year & $\bar\pi=5\%$ a year\\\midrule
\multicolumn{3}{@{}l}{\emph{Cross-type common mode $p^{\rm cc}$ per opportunity (human miss rate 0.05)}}\\
0 & 1H + 2AI, \$0.80m (0.080\%) & 1H + 1AI, \$0.75m (0.075\%) \\
0.0001 & 1H + 2AI, \$0.80m (0.080\%) & 1H + 1AI, \$0.75m (0.075\%) \\
0.0002 & 2H + 1AI, \$1.25m (0.125\%) & 1H + 1AI, \$0.75m (0.075\%) \\
0.0005 & none admissible & 1H + 1AI, \$0.75m (0.075\%) \\
0.0008 & none admissible & 1H + 1AI, \$0.75m (0.075\%) \\
0.001 & none admissible & 1H + 2AI, \$0.80m (0.080\%) \\
0.0015 & none admissible & 1H + 2AI, \$0.80m (0.080\%) \\
0.002 & none admissible & 2H + 1AI, \$1.25m (0.125\%) \\
0.0025 & none admissible & none admissible \\
\addlinespace\multicolumn{3}{@{}l}{\emph{Human miss rate on items the AI layers missed ($p^{\rm cc}=0$)}}\\
0.05 & 1H + 2AI, \$0.80m (0.080\%) & 1H + 1AI, \$0.75m (0.075\%) \\
0.1 & 2H + 1AI, \$1.25m (0.125\%) & 1H + 2AI, \$0.80m (0.080\%) \\
0.2 & 2H + 2AI, \$1.30m (0.130\%) & 1H + 2AI, \$0.80m (0.080\%) \\
0.5 & 4H + 2AI, \$2.30m (0.230\%) & 2H + 2AI, \$1.30m (0.130\%) \\
\bottomrule\end{tabular}
\par\vspace{2pt}\footnotesize\raggedright Notes: 24 opportunities a year at \$1 billion of revenue; human layer: miss rate 0.05, common mode 0.002, \$500{,}000; AI layer: miss rate 0.02, common mode 0.005, \$50{,}000, plus the ITGC and revalidation line of one in-scope AI application, \$204{,}672, for any design with AI layers; designs of up to six layers of each type. H: human layers; AI: AI layers; cost in millions of dollars and as a share of revenue. A design is admissible if the probability of at least one undetected material misstatement in the year is at most $\bar\pi$. The threshold is a scenario value; the table conditions on it and on the illustrative layer parameters, and it is not a control-design recommendation.
\end{table}}{}

\begin{table}[p]
\centering
\caption{Full-cost floors at full automation, percent of revenue unless stated}\label{tab:fullcost}
\footnotesize
\setlength{\tabcolsep}{3pt}
\begin{tabular}{@{}p{6.4cm}ccc p{5.6cm}@{}}
\toprule
Component & Fav. & Central & Cons. & Grounding and source\\
\midrule
% BEGIN ROWS GENERATED BY code/calibrate_horizons.py (tab:fullcost); do not edit by hand
\multicolumn{5}{@{}l}{\emph{Security, $c_{\mathrm{sec}}$}}\\
Security spending today, $c^{\rm sp}_0$ & 0.69 & 0.69 & 0.69 & (a) survey of security officers, flagged\\
Control effectiveness $L_{\rm sec}$ & 4 & 2 & 1 & (d) controls remove 98, 86, 63\% of loss\\
Exposure multiplier at full automation $\xi_{\rm exp}$ & 1 & 2 & 3 & (d) draft, write and transact rights\\
Risk cost $(c^{\rm sp}_0/L_{\rm sec})(1+L_{\rm sec}+\ln\xi_{\rm exp})$ & 0.86 & 1.27 & 2.14 & Lemma~\ref{lem:fullcost}(b)\\
Post-quantum migration & 0.04 & 0.09 & 0.24 & (d) 0.75, 1.5, 3\% of IT spend a year; IT spend $c^{\rm sp}_0$ over the security share (a); windows in \citet{moody2026transition}; only public-sector estimates exist \citep{omb2024report}\\
AI compliance at \$1 billion of revenue & 0.10 & 0.20 & 0.40 & vendor estimates for security attestation and ISO/IEC 42001, flagged; AI risk-framework staff, red-teaming and privacy assessments for 48 processes (d)\\
$c_{\mathrm{sec}}$ & 1.00 & 1.56 & 2.78 & \\
\multicolumn{5}{@{}l}{\emph{Maintenance and revalidation, $c_{\mathrm{mnt}}$}}\\
Upkeep $\delta_M$ (share of builds a year) & 0.10 & 0.17 & 0.40 & (b) \citet{svanberg2024beyond}; software support fees\\
Revalidation cost per build $\varrho_{\rm rv}$ & 0.05 & 0.10 & 0.20 & (d)\\
Model-upgrade interval $T_{\mathrm{up}}$ (years) & 1.75 & 1.22 & 0.99 & (c) vendor retirement dates\\
Builds $c_{\rm bld}$ & 1.2 & 2.4 & 4.8 & 48 processes at \$0.5 million (b), at \$1 billion of revenue (d)\\
$c_{\mathrm{mnt}}=(\delta_M+\varrho_{\rm rv}/T_{\mathrm{up}})c_{\rm bld}$ & 0.15 & 0.60 & 2.89 & Lemma~\ref{lem:fullcost}(a)\\
\multicolumn{5}{@{}l}{\emph{AI failure, $c_{\mathrm{err}}$}}\\
Error rate with no review $\bar u=\varepsilon^{\rm o}+\varepsilon^{\rm h}$ (per unit) & 0.050 & 0.085 & 0.085 & (d) scenario value: an omission rate per source sentence plus a hallucination rate per output sentence (favorable) or per summary or answer (central, conservative), which have different denominators, from source-grounded summarization and retrieval-based legal research; no review at full automation, Lemma~\ref{lem:fullcost}(c)\\
Loss per undetected error $L_{\rm err}$ (multiple of the unit's cost) & 0.25 & 0.5 & 1 & (d)\\
Errors and hallucination, $W_A\bar uL_{\rm err}$, with $W_A=60\%$ (d) & 0.74 & 2.54 & 5.07 & Lemma~\ref{lem:fullcost}(c)\\
Misbehavior, $W_A\varepsilon^{\rm m}(1-d^{\rm m})L_{\rm mis}\mathbb E\min\{Y,\bar Y\}$ & 0.03 & 0.67 & 3.01 & $\varepsilon^{\rm m}$ 0.05, 0.1, 0.2\% per task and $L_{\rm mis}=10$, $\bar Y=100$ (d); recall $d^{\rm m}$ 0.95, 0.6, 0.42 (a); tail index $p_Y$ 2.3, 1.5, 1.15 (b); Lemma~\ref{lem:fullcost}(d)\\
Common-mode failures, $\nu_{\rm cm}T_{\rm con}[W_{\rm cm}+H_{\rm cm}W_{\rm cm}^2/(2Y_{\rm fb})]$ & 0.16 & 0.66 & 1.97 & 0.5, 2, 6 day-long provider incidents a year (b); $W_{\rm cm}=W_A$, $W_{\rm cm}/(2Y_{\rm fb})=1$, one provider (d); Lemma~\ref{lem:fullcost}(e)\\
$c_{\mathrm{err}}$ & 0.93 & 3.86 & 10.05 & \\
\multicolumn{5}{@{}l}{\emph{Other, $c_{\mathrm{anc}}$}}\\
Management of AI spend (share of $c_c$) & 1.5 & 3 & 6 & (d)\\
$c_{\mathrm{anc}}$, with fixed incident-response and legal costs & 0.05 & 0.19 & 0.74 & (d); high-compute case 1.21\\
\midrule
$c_{\mathrm{full}}$ & 2.14 & 6.23 & 16.47 & high-compute case 7.25\\
% END ROWS GENERATED
\bottomrule
\end{tabular}

\smallskip
\begin{minipage}{0.97\linewidth}\footnotesize
Notes: the central column is the central case of Table~\ref{tab:ceiling}, and the high-compute case uses it with $c_c=0.40$. Financial misstatements are governed by the control standard of \autoref{prop:amin} and priced in $a_{\min}$, not in $c_{\mathrm{err}}$. The common-mode row treats $H_{\rm cm}$ as a lower bound on effective concentration. The AI-compliance row covers voluntary attestations and governance work; it excludes the obligations and penalties of the EU AI Act for high-risk systems, which apply from 2 December 2027 (Annex III systems) and 2 August 2028 (Annex I) \citep{eu2026reg1744}, and would rise for a firm that automates such uses. Every case of the table rests on scenario (class d) values and must be read conditionally. Values computed by \texttt{code/verify\_theory.py} and \texttt{code/calibrate\_horizons.py}.
\end{minipage}
\end{table}

\begin{remark}[Owning the compute]\label{rem:owncompute}
Let the compute behind $c_c$ cost $\mathrm{TCO}$ per dollar of revenue a year at cost, measured on a cash basis: the running cost of the capacity plus straight-line replacement of the assets that provide it. Let a vendor charge $(1+\upsilon_c)\mathrm{TCO}$. A firm that owns some of these assets holds capital $k_c=\kappa\,\mathrm{TCO}$, spends $s\,\mathrm{TCO}$ a year replacing them and leaves the rest of the cost, $(1-s)\mathrm{TCO}$, in $c_c$; for servers alone, with life $L_{\rm srv}$ and replacement share $s_{\rm srv}$, $s=s_{\rm srv}$ and $\kappa=s_{\rm srv}L_{\rm srv}$. Growth capital and replacement are the same equipment and take the same deduction rate $\psi$ (\autoref{rem:taxes}). By \autoref{cor:ceiling}, owning changes the ceiling of the Rule of 40 by $(1-\tau)\upsilon_c\mathrm{TCO}-(1-\psi)\tau s\,\mathrm{TCO}-(1-\psi\tau)\kappa\,\mathrm{TCO}\,G$. Under current U.S.\ law domestic servers take the 100\% allowance, $\psi=1$, and owning raises the cash ceiling if and only if $\upsilon_c>\kappa G$, whatever the tax rate.

The values below use this cash basis throughout. We take the assets and their lives from Epoch AI's cost model of a one-gigawatt data center \citep{michael2026servers}: servers costing \OwnEpochKSrv, the facility \OwnEpochKFac{} and the network \OwnEpochKNet{} billion dollars up front, running costs of \OwnEpochOpex{} billion a year, and lives of 5 years for IT equipment and 14 for the facility. On the cash basis the capacity costs \OwnTcoCash{} billion dollars a year, of which server replacement is \OwnSrvShare\%. A firm that owns only the servers, the case we take as central, has $\kappa=\OwnKapSrv$, and the markup must exceed \OwnMarkupFive\% at 5\% market growth and \OwnMarkupFifteen\% at 15\%. A firm that owns the whole data center, facility, network and land included, has $\kappa=\OwnKapAll$, which raises the bar to \OwnMarkupFiveAll\% and \OwnMarkupFifteenAll\% (with every asset at $\psi=1$; a building deducted more slowly raises it further). The source's headline share of 60\% is not on this basis: it is the servers' share of an annualized cost of \OwnEpochAnnTotal{} billion dollars a year whose capital charge includes a return on capital (the servers' annual charge is \OwnEpochCrf\% of their cost, a five-year annuity at about \OwnEpochRate\%, against \OwnEpochSl\% for straight-line replacement), which the cash score does not charge; a vendor that priced the capacity at that annualized cost would be charging a markup of \OwnEpochMarkup\% over its cash cost. With Meta's 5.5-year book life for IT equipment in place of Epoch's 5 years (Table~\ref{tab:param-sources}), the bars are \OwnMarkupFiveMeta{} and \OwnMarkupFifteenMeta\% for the servers alone and \OwnMarkupFiveAllMeta{} and \OwnMarkupFifteenAllMeta\% for the whole data center. A firm that elects out of the allowance, or whose servers are used predominantly abroad, has $\psi<1$ and a higher bar; at the lower bound $\psi=0$ and $\tau=0.18$ the bar for the servers alone is \OwnMarkupFiveNone\% and \OwnMarkupFifteenNone\%. Economic obsolescence faster than the assumed life raises the bar, and a test of value rather than of the steady-state score charges the capital at $r+\delta_K$ in place of the growth rate \citep{murthy2026transformation}. The comparison assumes that owned capacity is fully used. If it is used a share $u<1$ of the time, its capital cost per unit used is $1/u$ times that at full use, so the markup at which owning pays rises as utilization falls; operating staff and power limits, which the remark omits, raise it further.
\end{remark}
